\documentclass[10pt,a4paper]{amsart}
\usepackage{amsmath,amssymb,amsthm,graphicx}
% Retained preamble for the Stacks Project chapter snapshot in this directory.
% The chapter's own preamble.tex is a copy of the shared Stacks Project preamble; the block
% below is sliced MECHANICALLY, line for line, from it. Nothing was added, removed or
% redefined; the wrapper class, the xr-hyper/hyperref wiring and the externaldocument list
% are omitted because the assembler supplies the document class.
%
% lines taken from preamble.tex: 3, 5, 9, 11, 12, 156, 157, 158, ... (51 lines total)
%
% package coverage restored from the chapter's own preamble (2026-10-13): xy, xy-2cell

\usepackage[all]{xy}
\xyoption{2cell}
\UseAllTwocells

\newenvironment{reference}{\comment}{\endcomment}
\newenvironment{slogan}{\comment}{\endcomment}
\newenvironment{history}{\comment}{\endcomment}
\theoremstyle{plain}
\newtheorem{theorem}[subsection]{Theorem}
\newtheorem{proposition}[subsection]{Proposition}
\newtheorem{lemma}[subsection]{Lemma}
\theoremstyle{definition}
\newtheorem{definition}[subsection]{Definition}
\newtheorem{example}[subsection]{Example}
\newtheorem{exercise}[subsection]{Exercise}
\newtheorem{situation}[subsection]{Situation}
\theoremstyle{remark}
\newtheorem{remark}[subsection]{Remark}
\newtheorem{remarks}[subsection]{Remarks}
\numberwithin{equation}{subsection}
\def\lim{\mathop{\mathrm{lim}}\nolimits}
\def\colim{\mathop{\mathrm{colim}}\nolimits}
\def\Spec{\mathop{\mathrm{Spec}}}
\def\Hom{\mathop{\mathrm{Hom}}\nolimits}
\def\Ext{\mathop{\mathrm{Ext}}\nolimits}
\def\SheafHom{\mathop{\mathcal{H}\!\mathit{om}}\nolimits}
\def\SheafExt{\mathop{\mathcal{E}\!\mathit{xt}}\nolimits}
\def\Sch{\mathit{Sch}}
\def\Mor{\mathop{\mathrm{Mor}}\nolimits}
\def\Ob{\mathop{\mathrm{Ob}}\nolimits}
\def\Sh{\mathop{\mathit{Sh}}\nolimits}
\def\NL{\mathop{N\!L}\nolimits}
\def\CH{\mathop{\mathrm{CH}}\nolimits}
\def\proetale{{pro\text{-}\acute{e}tale}}
\def\etale{{\acute{e}tale}}
\def\QCoh{\mathit{QCoh}}
\def\Ker{\mathop{\mathrm{Ker}}}
\def\Im{\mathop{\mathrm{Im}}}
\def\Coker{\mathop{\mathrm{Coker}}}
\def\Coim{\mathop{\mathrm{Coim}}}
\DeclareMathSymbol{\boxtimes}{\mathbin}{AMSa}{"02}
\def\QCohstack{\mathcal{QC}\!\mathit{oh}}
\def\Cohstack{\mathcal{C}\!\mathit{oh}}
\def\Spacesstack{\mathcal{S}\!\mathit{paces}}
\def\Quotfunctor{\mathrm{Quot}}
\def\Hilbfunctor{\mathrm{Hilb}}
\def\Curvesstack{\mathcal{C}\!\mathit{urves}}
\def\Polarizedstack{\mathcal{P}\!\mathit{olarized}}
\def\Complexesstack{\mathcal{C}\!\mathit{omplexes}}
\def\Pic{\mathop{\mathrm{Pic}}\nolimits}
\def\Picardstack{\mathcal{P}\!\mathit{ic}}
\def\Picardfunctor{\mathrm{Pic}}
\def\Deformationcategory{\mathcal{D}\!\mathit{ef}}
\usepackage{cleveref}
\crefname{originallemma}{original lemma}{original lemmas}
\Crefname{originallemma}{Original lemma}{Original lemmas}
\makeatletter
\newlabel{algebra-lemma-lift-nth-roots}{{0CAP}{source}{}{sourcepointer-1}{}}
\newlabel{algebra-lemma-lift-nth-roots@cref}{{[originallemma][0][]0CAP}{[1][0][]source}{}{}{}}
\makeatother
\crefname{originallemma}{original lemma}{original lemmas}
\Crefname{originallemma}{Original lemma}{Original lemmas}
\makeatletter
\newlabel{coherent-lemma-flat-base-change-cohomology}{{02KH}{source}{}{sourcepointer-2}{}}
\newlabel{coherent-lemma-flat-base-change-cohomology@cref}{{[originallemma][0][]02KH}{[1][0][]source}{}{}{}}
\makeatother
\crefname{originallemma}{original lemma}{original lemmas}
\Crefname{originallemma}{Original lemma}{Original lemmas}
\makeatletter
\newlabel{descent-definition-descent-datum-quasi-coherent}{{023B}{source}{}{sourcepointer-3}{}}
\newlabel{descent-definition-descent-datum-quasi-coherent@cref}{{[originallemma][0][]023B}{[1][0][]source}{}{}{}}
\makeatother
\crefname{originallemma}{original lemma}{original lemmas}
\Crefname{originallemma}{Original lemma}{Original lemmas}
\makeatletter
\newlabel{descent-proposition-fpqc-descent-quasi-coherent}{{023T}{source}{}{sourcepointer-4}{}}
\newlabel{descent-proposition-fpqc-descent-quasi-coherent@cref}{{[originallemma][0][]023T}{[1][0][]source}{}{}{}}
\makeatother
\crefname{originallemma}{original lemma}{original lemmas}
\Crefname{originallemma}{Original lemma}{Original lemmas}
\makeatletter
\newlabel{more-morphisms-lemma-torsion-pic-thickening}{{0C6S}{source}{}{sourcepointer-5}{}}
\newlabel{more-morphisms-lemma-torsion-pic-thickening@cref}{{[originallemma][0][]0C6S}{[1][0][]source}{}{}{}}
\makeatother
\crefname{originallemma}{original lemma}{original lemmas}
\Crefname{originallemma}{Original lemma}{Original lemmas}
\makeatletter
\newlabel{varieties-lemma-change-fields-pic}{{0CC5}{source}{}{sourcepointer-6}{}}
\newlabel{varieties-lemma-change-fields-pic@cref}{{[originallemma][0][]0CC5}{[1][0][]source}{}{}{}}
\makeatother
\title{Torsion line bundles descend along purely inseparable extensions}
\author{The Stacks Project Authors}
\date{stacks-project:pic.tex:local-snapshot-20261012; GNU FDL 1.2 or later}
\begin{document}
\maketitle
\noindent\textit{Editorial note.} Statement span is lines 1625-1635 of sources/stacks-pic/pinned-original-source.tex; proof span is lines 1637-1718, the complete author proof (80 body lines).\par\medskip
\noindent\textit{Editorial note.} Navigation stubs: 6 cited-but-not-included labels are supplied as non-load-bearing pointers with their pinned tags where the snapshot tag index has them.\par\medskip
\noindent\textit{Editorial note.} extra\_packages is empty on purpose: the retained preamble already loads xy with options and a second declaration clashes in the assembler wrapper.\par\medskip
\noindent\textit{Editorial note.} Text is retained verbatim, including the author's own typos and placeholders.\par\medskip
\noindent\textit{Editorial note.} Queue index 3 of queue\_pic.json (pb=3551, stmt 1625-1635, proof 1637-1718).\par\medskip
\noindent\textit{Editorial note.} stacks-pic wave-1 (pi01): fresh adjudication accept (2026-10-14); pool per standing queue hypercovering -> pic (SMC reassigned to Zhen per Chief 08:5x); pic.tex 27 labels/23 with proof judged 21 accept / 2 reject / 4 pb=0 env-dup rows excluded; spans cite-free and reference/slogan/history-free so no patchenv needed.\par\medskip
\noindent\textit{Allowed standard prerequisites (editorial).}\begin{itemize}
\item Explicit prerequisite boundary. The cited inputs are prerequisites; the argument in the retained span is not replaced by any cited result.
\end{itemize}\medskip
\ifdefined\hypertarget\hypertarget{sourcepointer-1}{}\fi
\noindent\textit{Source reference.} Original lemma 0CAP: sources/stacks-pic/pinned-original-source.tex, source lines 1718--1718. Cited by the retained statement or proof but not part of the collected spans; supplied only as a navigation pointer. The argument does not depend on its content.\par\medskip
\ifdefined\hypertarget\hypertarget{sourcepointer-2}{}\fi
\noindent\textit{Source reference.} Original lemma 02KH: sources/stacks-pic/pinned-original-source.tex, source lines 1718--1718. Cited by the retained statement or proof but not part of the collected spans; supplied only as a navigation pointer. The argument does not depend on its content.\par\medskip
\ifdefined\hypertarget\hypertarget{sourcepointer-3}{}\fi
\noindent\textit{Source reference.} Original lemma 023B: sources/stacks-pic/pinned-original-source.tex, source lines 1718--1718. Cited by the retained statement or proof but not part of the collected spans; supplied only as a navigation pointer. The argument does not depend on its content.\par\medskip
\ifdefined\hypertarget\hypertarget{sourcepointer-4}{}\fi
\noindent\textit{Source reference.} Original lemma 023T: sources/stacks-pic/pinned-original-source.tex, source lines 1718--1718. Cited by the retained statement or proof but not part of the collected spans; supplied only as a navigation pointer. The argument does not depend on its content.\par\medskip
\ifdefined\hypertarget\hypertarget{sourcepointer-5}{}\fi
\noindent\textit{Source reference.} Original lemma 0C6S: sources/stacks-pic/pinned-original-source.tex, source lines 1718--1718. Cited by the retained statement or proof but not part of the collected spans; supplied only as a navigation pointer. The argument does not depend on its content.\par\medskip
\ifdefined\hypertarget\hypertarget{sourcepointer-6}{}\fi
\noindent\textit{Source reference.} Original lemma 0CC5: sources/stacks-pic/pinned-original-source.tex, source lines 1718--1718. Cited by the retained statement or proof but not part of the collected spans; supplied only as a navigation pointer. The argument does not depend on its content.\par\medskip
\begin{lemma}
\label{lemma-torsion-descends}
Let $k$ be a field of characteristic $p > 0$. Let $X$ be a
quasi-compact and quasi-separated scheme over $k$ with
$H^0(X, \mathcal{O}_X) = k$. Let $n$ be an integer prime to $p$.
Then the map
$$
\Pic(X)[n] \longrightarrow \Pic(X_{k'})[n]
$$
is bijective for any purely inseparable extension $k'/k$.
\end{lemma}

\begin{proof}
First we observe that the map $\Pic(X) \to \Pic(X_{k'})$
is injective by Varieties, Lemma \ref{varieties-lemma-change-fields-pic}.
Hence we have to show the map in the lemma is surjective.
Let $\mathcal{L}$ be an invertible $\mathcal{O}_{X_{k'}}$-module
which has order dividing $n$ in $\Pic(X_{k'})$.
Choose an isomorphism
$\alpha : \mathcal{L}^{\otimes n} \to \mathcal{O}_{X_{k'}}$
of invertible modules.
We will prove that we can descend the pair $(\mathcal{L}, \alpha)$
to $X$.

\medskip\noindent
Set $A = k' \otimes_k k'$. Since $k'/k$ is purely inseparable,
the kernel of the multiplication map $A \to k'$ is a
locally nilpotent ideal $I$ of $A$. Observe that
$$
X_A = X \times_{\Spec(k)} \Spec(A) = X_{k'} \times_X X_{k'}
$$
comes with two projections $\text{pr}_i : X_A \to X_{k'}$, $i = 0, 1$
which agree over $A/I$. Hence the invertible modules
$\mathcal{L}_i = \text{pr}_i^*\mathcal{L}$
agree over the closed subscheme $X_{A/I} = X_{k'}$.
Since $X_{A/I} \to X_A$ is a thickening and since
$\mathcal{L}_i$ are $n$-torsion, we see that there
exists an isomorphism $\varphi : \mathcal{L}_0 \to \mathcal{L}_1$ by
More on Morphisms, Lemma \ref{more-morphisms-lemma-torsion-pic-thickening}.
We may pick $\varphi$ to reduce to the identity modulo $I$. Namely,
$H^0(X, \mathcal{O}_X) = k$ implies
$H^0(X_{k'}, \mathcal{O}_{X_{k'}}) = k'$ by
Cohomology of Schemes, Lemma \ref{coherent-lemma-flat-base-change-cohomology}
and $A \to k'$ is surjective hence we can adjust $\varphi$ by multiplying by
a suitable element of $A$. Consider the map
$$
\lambda : \mathcal{O}_{X_A}
\xrightarrow{\text{pr}_0^*\alpha^{-1}}
\mathcal{L}_0^{\otimes n}
\xrightarrow{\varphi^{\otimes n}}
\mathcal{L}_1^{\otimes n}
\xrightarrow{\text{pr}_0^*\alpha}
\mathcal{O}_{X_A}
$$
We can view $\lambda$ as an element of $A$ because
$H^0(X_A, \mathcal{O}_{X_A}) = A$ (same reference as above).
Since $\varphi$ reduces to the identity modulo $I$ we
see that $\lambda = 1 \bmod I$. Then there is a unique
$n$th root of $\lambda$ in $1 + I$
(Algebra, Lemma \ref{algebra-lemma-lift-nth-roots})
and after multiplying $\varphi$ by its inverse we get $\lambda = 1$.
We claim that $(\mathcal{L}, \varphi)$ is a descent datum
for the fpqc covering $\{X_{k'} \to X\}$
(Descent, Definition \ref{descent-definition-descent-datum-quasi-coherent}).
If true, then $\mathcal{L}$ is the pullback of an invertible
$\mathcal{O}_X$-module $\mathcal{N}$ by
Descent, Proposition \ref{descent-proposition-fpqc-descent-quasi-coherent}.
Injectivity of the map on Picard groups shows that $\mathcal{N}$
is a torsion element of $\Pic(X)$ of the same order as $\mathcal{L}$.

\medskip\noindent
Proof of the claim. To see this we have to verify that
$$
\text{pr}_{12}^*\varphi \circ
\text{pr}_{01}^*\varphi =
\text{pr}_{02}^*\varphi
\quad\text{on}\quad
X_{k'} \times_X X_{k'} \times_X X_{k'} = X_{k' \otimes_k k' \otimes_k k'}
$$
As before the diagonal morphism
$\Delta : X_{k'} \to X_{k' \otimes_k k' \otimes_k k'}$
is a thickening. The left and right hand sides of the equality signs
are maps $a, b : p_0^*\mathcal{L} \to p_2^*\mathcal{L}$ compatible with
$p_0^*\alpha$ and $p_2^*\alpha$ where
$p_i : X_{k' \otimes_k k' \otimes_k k'} \to X_{k'}$
are the projection morphisms. Finally, $a, b$ pull back to the
same map under $\Delta$.
Affine locally (in local trivializations) this means that
$a, b$ are given by multiplication by invertible functions
which reduce to the same function modulo a locally nilpotent
ideal and which have the same $n$th powers. Then it follows
from Algebra, Lemma \ref{algebra-lemma-lift-nth-roots}
that these functions are the same.
\end{proof}

\end{document}
